1: % Reference entries, S (AN-11).  Law 1: os/iword/app/face_a.c F10 (the
2: % offered name, a typed character replaces it; save_named: an existing file
3: % asks "Replace A:\NAME? No (Yes)", y replaces, any other key asks for a
4: % name again; write_file: flatten first, "ERROR: File I/O error", "ERROR:
5: % Disk Full--Press Any Key to Continue"), wp5.c iw_wp5_write (the 327-byte
6: % prefix of WordPerfect 5.1 and the body), M_SIZE (IW_SIZE_MENU,
7: % IW_SIZE_ATTRS), paint_status (the status line: the name, Block on, Doc Pg
8: % Ln Pos, POS with Caps Lock; iw_fmt_inches).  Gated byte for byte against
9: % the real WordPerfect 5.1 by test-iword-face and test-flair-iword (MEMO ==
10: % MINT02's prefix + the body).
11: % Search (F2, Shift-F2) and Setup (Shift-F1) are P3 (bead initech-qvgd,
12: % WL-0146): face_a.c function_key (IW_FK_search_forward / _backward,
13: % IW_FK_setup from spec/iword/fkeys.json), search_key (the string kept for
14: % the next Search and replaced by the first key typed; F2 / Shift-F2 / Esc
15: % search, F1 cancels with the cursor kept, Up / Down the direction, Left /
16: % Right / Backspace / Delete by whole codes, F6 / F8 an attribute's on or
17: % off code by the state at that point of the string, Tab [Tab], F4
18: % [->Indent], Enter [HRt], Ctrl-Enter [HPg]; 256 bytes, QUERY_CAP),
19: % os/iword/core/search.c iw_search (a lower-case letter finds either case,
20: % a capital only a capital; the cursor after the match forward, at its
21: % start backward; "* Not found *"), paint_setup / setup_key / settings_save
22: % (the Setup screens, the grey rows, Environment 3, Backup 1 / 2, Units 8
23: % with i c p w u; F7 one level up, F1 / Esc to the document; IWORD.SET,
24: % spec/iword/settings.h, eight bytes; the units and the backup choice are
25: % recorded and shown, nothing else reads them).
26: % GATES: test-pubs-iword test-iword-face test-flair-iword test-iword-wp5-roundtrip test-iword-codes test-iword-reveal-codes test-iword-search test-iword-setup test-flair-iword-search test-flair-iword-setup
27: \Gates{test-pubs-iword test-iword-face test-flair-iword test-iword-wp5-roundtrip test-iword-codes test-iword-reveal-codes test-iword-search test-iword-setup test-flair-iword-search test-flair-iword-setup}
28: \begin{Feature}{Save}
29: Save writes the document on the screen to the data volume as a file, and
30: leaves it on the screen so that you can go on working.  Save often: a
31: document that is only on the screen is lost if the computer is switched off.
32: 
33: To save a document,
34: \begin{WPSteps}
35: \item Press \FK{Save}{F10}.  The status line shows \Scr{Document to be
36: Saved:} and, if the document has a name, that name.
37: \PullDown{Select \Mi{S}ave from the \Mi{F}ile menu.}
38: \item To save under the name shown, press \K{Enter}.
39: \Or
40: Type a new name (the first character you type replaces the name shown) and
41: press \K{Enter}.
42: \item If a document with that name is already on the data volume, the status
43: line asks \Scr{Replace A:\wbs\textit{name}? No (Yes)}.  Type \Typed{y} to
44: replace it with the document on the screen.  Any other key asks you for a
45: name again, so that you can choose another.
46: \end{WPSteps}
47: \Documents{wpkey}{F10}
48: When the document is saved, its name appears at the left of the status line.
49: 
50: \BeginReplay{ref-save}
51: For example, type \TypeRun{Budget memo}, press \Press{Save}{F10}, type
52: \TypeName{budget} and press \PressKey{Enter}.  Press \Press{Save}{F10}
53: again: InitechWord offers the name.
54: \WordScreen[rows=1]{ref-save-1}{Save offers the document's name}
55: Press \PressKey{Enter}.  The file \Scr{BUDGET} exists, so InitechWord asks:
56: \WordScreen[rows=1]{ref-save-2}{Saving over a document that is already on the data volume}
57: Type \TypeName{y}: the document is saved again.
58: \EndReplay
59: 
60: \Notes
61: \Topic{File Format}
62: InitechWord saves documents in the format of WordPerfect 5.1, writing the
63: text and codes as WordPerfect 5.1 does (see \InRef{WordPerfect 5.1
64: Documents} and Appendix D).
65: 
66: \Topic{Messages}
67: \Scr{ERROR: Invalid file name} means the name breaks the rules for names
68: (see \InRef{Filenames}).  \Scr{ERROR: File I/O error} means the file could not
69: be made on the data volume.  \Scr{ERROR: Disk Full--Press Any Key to
70: Continue} means the data volume has no room for the document; the document
71: is still on the screen.  After each of these, press any key to return to
72: the name prompt.  \Scr{ERROR: Document too large} means the document is
73: larger than a file InitechWord can write (see Appendix H); press any key to
74: return to the document.
75: 
76: \Topic{Saving When You Exit}
77: Exit (F7) asks whether to save the document before you leave it, and saves
78: it as Save does (see \InRef{Exit}).
79: 
80: \SeeAlso{Exit; Filenames; Retrieve}
81: \end{Feature}
82: 
83: \begin{Feature}{Screen}
84: The editing screen is the InitechWord window as you type.  It shows 23
85: lines of the document, 76 characters across, and the status line under
86: them.  The first column of the screen is the left margin of a new document,
87: one inch from the left edge of the paper.
88: 
89: \Notes
90: \Topic{Colors}
91: Text is white on a blue field.  The cursor is the cell drawn in reverse,
92: blue on white.  A block is black on grey; so is the ruler of Reveal Codes.
93: 
94: \Topic{Faces}
95: Bold, italic and underlined text is drawn in that face.  Text with any other
96: appearance or size is drawn as ordinary text (see \InRef{Appearance} and
97: \InRef{Size}).
98: 
99: \Topic{Hard Pages and Spacing}
100: A hard page is drawn as a row of equal signs below the line it ends.  With
101: line spacing of more than one, blank rows are drawn under each line (see
102: \InRef{Line Spacing}).
103: 
104: \Topic{Margins Under One Inch}
105: The first column of the screen is always one inch from the edge of the
106: paper.  With a left margin of less than one inch, the start of each line
107: lies to the left of the window and is not shown; the status line's Pos still
108: tells where the cursor is.
109: 
110: \Topic{Scrolling}
111: When the cursor moves above the top line or below the bottom line, the
112: screen moves to keep it in view.
113: 
114: \SeeAlso{Reveal Codes; Status Line; Window}
115: \end{Feature}
116: 
117: \begin{Feature}{Search}
118: Search finds text, codes or both in the document.  It searches forward,
119: from the cursor toward the end of the document, or backward, from the
120: cursor toward the beginning, and stops at the first place where it finds
121: what you typed.
122: 
123: To search,
124: \begin{WPSteps}
125: \item Press \FK{\rarr Search}{F2} to search forward
126: \Or
127: Press \FK{\larr Search}{Shift-F2} to search backward.
128: \item The status line shows \Scr{-> Srch:} (forward) or \Scr{<- Srch:}
129: (backward), followed by what you last searched for, if anything.  Type the
130: text you want to find.  The first character you type replaces the text
131: shown.
132: \item Press \K{F2} to start the search.
133: \end{WPSteps}
134: \FlairShot{search}{640}{440}{The search string typed at the \Scr{-> Srch:} prompt}
135: \Documents{wpkey}{F2}
136: \Documents{wpkey}{Shift-F2}
137: \Documents{wpkey}{Srch:F2}
138: \Documents{wpkey}{Srch:Esc}
139: If InitechWord finds the text, the cursor stops just after it when you
140: search forward, and on its first character when you search backward.  If it
141: does not, the status line shows \Scr{* Not found *} and the cursor stays
142: where it was; press any key to go on.
143: 
144: \BeginReplay{ref-search}
145: For example, type
146: \Type{Every report needs a cover sheet.  The cover sheets are in the supply room.}
147: press \PressKey{Enter} and type
148: \Type{Ask Milton for Cover Sheets if the room is empty.}
149: Press \PressKey{Home, Home, Up} to go to the top of the document.  Press
150: \Press{\rarr Search}{F2} and type \TypeName{cover sheet}:
151: \WordScreen[rows=3]{ref-search-1}{The search string on the status line}
152: Press \PressKey{F2}.  The cursor stops after the first \textit{cover sheet}:
153: \WordScreen[rows=3]{ref-search-2}{The cursor after the text found}
154: Press \PressKey{F2} twice.  The first F2 shows \Scr{-> Srch:cover sheet}
155: again; the second searches for it, and the cursor stops after the
156: \textit{cover sheet} of \textit{cover sheets}.  Press \PressKey{F2} twice
157: more:\Keys{F2 F2}
158: \WordScreen[rows=3]{ref-search-3}{A string in lower case finds capitals as well}
159: Search found \textit{Cover Sheet}: a lower-case letter in the search string
160: finds the letter in either case.
161: \EndReplay
162: 
163: \Notes
164: \Topic{Block}
165: While Block is on, F2 and Shift-F2 do nothing.  Turn Block off to search.
166: 
167: \Topic{Canceling a Search}
168: Press \FK{Cancel}{F1} at the \Scr{Srch:} prompt to leave without searching.
169: The cursor stays where it was, and the search string is kept for the next
170: search.
171: \Documents{wpkey}{Srch:F1}
172: 
173: \Topic{Capital Letters}
174: A lower-case letter in the search string finds the same letter in upper or
175: lower case: \Typed{memo} finds \Scr{memo}, \Scr{Memo} and \Scr{MEMO}.  A
176: capital letter finds only a capital: \Typed{Memo} finds \Scr{Memo} and
177: \Scr{MEMO}, but not \Scr{memo}.
178: 
179: \Topic{Codes in the Search String}
180: To search for a code, press its key at the \Scr{Srch:} prompt.  The code is
181: put in the string as Reveal Codes shows it:
182: \par\addvspace{4pt}
183: \noindent\begin{tabular}{@{}ll@{}}
184: \textbf{Key} & \textbf{Code in the string}\\
185: \K{F6} & \Code{[BOLD]}, or \Code{[bold]} after a \Code{[BOLD]} not yet ended\\
186: \K{F8} & \Code{[UND]}, or \Code{[und]} after an \Code{[UND]} not yet ended\\
187: \K{Tab} & \Code{[Tab]}\\
188: \K{F4} & \Code{[\rarr Indent]}\\
189: \K{Enter} & \Code{[HRt]}\\
190: \K{Ctrl-Enter} & \Code{[HPg]}\\
191: \end{tabular}
192: \par\addvspace{4pt}
193: \Documents{wpkey}{Srch:F6}
194: \Documents{wpkey}{Srch:F8}
195: \Documents{wpkey}{Srch:Tab}
196: \Documents{wpkey}{Srch:F4}
197: \Documents{wpkey}{Srch:Enter}
198: \Documents{wpkey}{Srch:Ctrl-Enter}
199: A string can hold codes and text together: \Code{[BOLD]}\Typed{Note} finds
200: the word \textit{Note} only where bold begins just before it.  Codes of other
201: kinds cannot be put in the search string.
202: 
203: \BeginReplay{ref-search-code}
204: For example, type \TypeRun{Read the }, press \Press{Bold}{F6}, type
205: \TypeRun{first}, press \Press{Bold}{F6} again and type
206: \TypeRun{ page.  The first page is the summary.}  Press
207: \PressKey{Home, Home, Up}, press \Press{\rarr Search}{F2}, press
208: \PressKey{F6} and type \TypeName{first}:
209: \WordScreen[rows=1]{ref-search-code-1}{A code in the search string}
210: Press \PressKey{F2}.  The cursor stops after the bold \textit{first}, not
211: after the second one, which has no code before it.
212: \WordScreen[rows=1]{ref-search-code-2}{The bold word found}
213: \EndReplay
214: 
215: \Topic{Direction}
216: At the \Scr{Srch:} prompt, \K{Up} changes the search to backward
217: (\Scr{<- Srch:}) and \K{Down} to forward (\Scr{-> Srch:}).  \K{F2},
218: \K{Shift-F2} and \K{Esc} all start the search in the direction the prompt
219: shows.
220: \Documents{wpkey}{Srch:Up}
221: \Documents{wpkey}{Srch:Down}
222: 
223: \Topic{Editing the Search String}
224: \K{Left} and \K{Right} move the cursor through the search string, and
225: \K{Backspace} and \K{Delete} erase, as on the editing screen; each moves
226: past or erases a whole code at a time.  Once you have pressed one of them,
227: typing no longer replaces the string: what you type goes in at the cursor.
228: \Documents{wpkey}{Srch:Left}
229: \Documents{wpkey}{Srch:Right}
230: \Documents{wpkey}{Srch:Backspace}
231: \Documents{wpkey}{Srch:Delete}
232: 
233: \Topic{Keys That Act as Other Keys}
234: At the \Scr{Srch:} prompt InitechWord does not look at Shift, Alt or Ctrl on
235: these keys: each combination does what the key alone does.
236: \par\addvspace{4pt}
237: \noindent{\small\begin{tabular}{@{}lll@{}}
238: \WPPromptAlias{Srch:}{Shift-F1}{F1}, & \WPPromptAlias{Srch:}{Alt-F1}{F1}, & \WPPromptAlias{Srch:}{Ctrl-F1}{F1}\\
239: \WPPromptAlias{Srch:}{Shift-F2}{F2}, & \WPPromptAlias{Srch:}{Alt-F2}{F2}, & \WPPromptAlias{Srch:}{Ctrl-F2}{F2}\\
240: \WPPromptAlias{Srch:}{Shift-Esc}{Esc}, & \WPPromptAlias{Srch:}{Alt-Esc}{Esc}, & \WPPromptAlias{Srch:}{Ctrl-Esc}{Esc}\\
241: \WPPromptAlias{Srch:}{Shift-Up}{Up}, & \WPPromptAlias{Srch:}{Alt-Up}{Up}, & \WPPromptAlias{Srch:}{Ctrl-Up}{Up}\\
242: \WPPromptAlias{Srch:}{Shift-Down}{Down}, & \WPPromptAlias{Srch:}{Alt-Down}{Down}, & \WPPromptAlias{Srch:}{Ctrl-Down}{Down}\\
243: \WPPromptAlias{Srch:}{Shift-Left}{Left}, & \WPPromptAlias{Srch:}{Alt-Left}{Left}, & \WPPromptAlias{Srch:}{Ctrl-Left}{Left}\\
244: \WPPromptAlias{Srch:}{Shift-Right}{Right}, & \WPPromptAlias{Srch:}{Alt-Right}{Right}, & \WPPromptAlias{Srch:}{Ctrl-Right}{Right}\\
245: \WPPromptAlias{Srch:}{Shift-Backspace}{Backspace}, & \WPPromptAlias{Srch:}{Alt-Backspace}{Backspace}, & \WPPromptAlias{Srch:}{Ctrl-Backspace}{Backspace}\\
246: \WPPromptAlias{Srch:}{Shift-Delete}{Delete}, & \WPPromptAlias{Srch:}{Alt-Delete}{Delete}, & \WPPromptAlias{Srch:}{Ctrl-Delete}{Delete}\\
247: \WPPromptAlias{Srch:}{Shift-Tab}{Tab}, & \WPPromptAlias{Srch:}{Alt-Tab}{Tab}, & \WPPromptAlias{Srch:}{Ctrl-Tab}{Tab}\\
248: \WPPromptAlias{Srch:}{Shift-Enter}{Enter}, & \WPPromptAlias{Srch:}{Alt-Enter}{Enter} & \\
249: \end{tabular}}
250: \par\addvspace{4pt}
251: (Ctrl-Enter is the hard page code, as above.)
252: 
253: \Topic{Length of the String}
254: The search string holds up to 256 bytes of the document file: one for each
255: character, \Code{[HRt]} or \Code{[HPg]}, and the length given in Appendix D
256: for each other code.  Keys typed after the string is full are ignored.
257: 
258: \Topic{Repeating a Search}
259: InitechWord keeps the search string until you leave InitechWord.  To find the next
260: place, press F2 twice: once to show the string, once to search.
261: 
262: \Topic{The Document Is Not Changed}
263: Searching moves only the cursor.  It does not change the document, and a
264: document that was saved is still saved after a search.
265: 
266: \Topic{Words Within Words}
267: Search finds the string wherever it occurs, also inside a longer word:
268: \Typed{form} finds \Scr{form}, \Scr{forms} and \Scr{platform}.  To find a
269: whole word, put a space before or after it in the string.
270: 
271: \SeeAlso{Block; Codes; Cursor Movement; Replace}
272: \end{Feature}
273: 
274: \begin{Feature}{Setup}
275: Setup records two preferences of InitechWord: the units in which
276: measurements are to be shown, and whether a copy of a document is to be kept
277: when it is saved over.  InitechWord keeps them in a file on the data volume,
278: so that they are the same each time you start it.
279: 
280: To change a setting,
281: \begin{WPSteps}
282: \item Press \FK{Setup}{Shift-F1}.  The Setup screen appears.
283: \item Select \Ch{Environment}{3}.
284: \item Select \Ch{Backup Options}{1} or \Ch{Units of Measure}{8}, then the
285: setting, and give its value (see the Notes below).
286: \item Press \FK{Exit}{F7} to go back one screen at a time, or \FK{Cancel}{F1}
287: to return to the document at once.
288: \end{WPSteps}
289: \Documents{wpkey}{Shift-F1}
290: \Documents{wpmenu}{Shift-F1,3=Environment}
291: \Documents{wpmenu}{Shift-F1,3,1=Backup Options}
292: \Documents{wpmenu}{Shift-F1,3,8=Units of Measure}
293: 
294: \FlairShot{setup}{640}{440}{The Setup screen}
295: The choices printed in grey are WordPerfect settings InitechWord does not
296: have; typing their numbers does nothing.
297: 
298: \BeginReplay{ref-setup}
299: For example, press \Press{Setup}{Shift-F1} and select
300: \Choose{Environment}{3} and \Choose{Units of Measure}{8}:
301: \WordScreen[rows=13]{ref-setup-2}{Setup: Units of Measure}
302: Press \FK{Exit}{F7}\Keys{F7} twice\Keys{F7} and \FK{Cancel}{F1}\Keys{F1} to
303: return to the document.
304: \EndReplay
305: 
306: \Notes
307: \Topic{Backup Options}
308: Select \Ch{Backup Options}{1} on the Environment screen.  Timed Document
309: Backup (1) is printed in grey: InitechWord has none.  Select
310: \Ch{Original Document Backup}{2}; the status line asks
311: \Scr{Original Document Backup? No (Yes)}.  Type \Typed{y} or \Typed{n}.
312: The screen shows the new value.
313: \Documents{wpmenu}{Shift-F1,3,1,2=Original Document Backup}
314: InitechWord records the answer, but in this release it keeps no backup copy
315: either way: Save replaces a document as described under \InRef{Save}.
316: 
317: \Topic{Exit and Cancel}
318: On every Setup screen \FK{Exit}{F7} goes back one screen: from Units of
319: Measure or Backup Options to Environment, from Environment to Setup, and
320: from Setup to the document.  \FK{Cancel}{F1} and \K{Esc} return to the
321: document from any Setup screen.
322: 
323: \Topic{The Setup File}
324: InitechWord keeps the settings in the file \Scr{IWORD.SET} in the top level
325: of the data volume, and writes it each time you change a setting.  The file
326: is eight bytes long: the letters \Scr{IWS1}, the letter of the unit for
327: margins and tabs, the letter of the unit for the status line, 1 if Original
328: Document Backup is Yes (otherwise 0), and a 0.  If the file is missing or
329: not in this form, InitechWord starts with Inches and No.  If the file cannot
330: be written, the status line shows \Scr{ERROR: File I/O error} and the setting
331: is not changed.
332: 
333: \Topic{Units of Measure}
334: Select \Ch{Units of Measure}{8} on the Environment screen, then
335: \Ch{Display and Entry of Numbers for Margins, Tabs, etc.}{1} or
336: \Ch{Status Line Display}{2}.  The status line shows
337: \Scr{Units of Measure:}; type the letter of a unit:
338: \par\addvspace{4pt}
339: \noindent\begin{tabular}{@{}ll@{}}
340: \textbf{Letter} & \textbf{Unit}\\
341: \Typed{i} & Inches\\
342: \Typed{c} & Centimeters\\
343: \Typed{p} & Points\\
344: \Typed{w} & 1200ths of an inch\\
345: \Typed{u} & WordPerfect 4.2 units\\
346: \end{tabular}
347: \par\addvspace{4pt}
348: \Documents{wpmenu}{Shift-F1,3,8,1=Display and Entry of Numbers for Margins, Tabs, etc.}
349: \Documents{wpmenu}{Shift-F1,3,8,1,i=Inches}
350: \Documents{wpmenu}{Shift-F1,3,8,1,c=Centimeters}
351: \Documents{wpmenu}{Shift-F1,3,8,1,p=Points}
352: \Documents{wpmenu}{Shift-F1,3,8,1,w=1200ths of an inch}
353: \Documents{wpmenu}{Shift-F1,3,8,1,u=WordPerfect 4.2 units}
354: \Documents{wpmenu}{Shift-F1,3,8,2=Status Line Display}
355: \Documents{wpmenu}{Shift-F1,3,8,2,i=Inches}
356: \Documents{wpmenu}{Shift-F1,3,8,2,c=Centimeters}
357: \Documents{wpmenu}{Shift-F1,3,8,2,p=Points}
358: \Documents{wpmenu}{Shift-F1,3,8,2,w=1200ths of an inch}
359: \Documents{wpmenu}{Shift-F1,3,8,2,u=WordPerfect 4.2 units}
360: Any other key is ignored.  The Units of Measure screen shows the unit
361: chosen.  In this release InitechWord records the units but goes on showing
362: and taking every measurement in inches (see \InRef{Units of Measure}).
363: 
364: \SeeAlso{Save; Units of Measure}
365: \end{Feature}
366: 
367: \begin{Feature}{Size}
368: Size changes the size of text: superscript and subscript, and the sizes from
369: fine to extra large.  Like the appearances, each size is a pair of codes
370: around the text.
371: 
372: To give text a size as you type it,
373: \begin{WPSteps}
374: \item Press \FK{Font}{Ctrl-F8}, then select \Ch{Size}{1}.  The Size menu
375: appears on the status line:
376: \StatusLine{1 Suprscpt; 2 Subscpt; 3 Fine; 4 Small; 5 Large; 6 Vry Large; 7 Ext Large: 0}
377: \item Select the size by its number.
378: \item Type the text.
379: \item Press \KC{Ctrl-F8}{1} and select the same size again to move past its
380: off code
381: \Or
382: Press \FK{Font}{Ctrl-F8} and select \Ch{Normal}{3}.
383: \end{WPSteps}
384: \Documents{wpmenu}{Ctrl-F8,1=Size}
385: \Documents{wpmenu}{Ctrl-F8,1,1=Suprscpt}
386: \Documents{wpmenu}{Ctrl-F8,1,2=Subscpt}
387: \Documents{wpmenu}{Ctrl-F8,1,3=Fine}
388: \Documents{wpmenu}{Ctrl-F8,1,4=Small}
389: \Documents{wpmenu}{Ctrl-F8,1,5=Large}
390: \Documents{wpmenu}{Ctrl-F8,1,6=Vry Large}
391: \Documents{wpmenu}{Ctrl-F8,1,7=Ext Large}
392: \par\addvspace{6pt}
393: \noindent\begin{tabular}{@{}lll@{}}
394: \textbf{Selection} & \textbf{Size} & \textbf{Codes}\\
395: Suprscpt (1) & superscript & \Code{[SUPRSCPT]} \Code{[suprscpt]}\\
396: Subscpt (2) & subscript & \Code{[SUBSCPT]} \Code{[subscpt]}\\
397: Fine (3) & fine & \Code{[FINE]} \Code{[fine]}\\
398: Small (4) & small & \Code{[SMALL]} \Code{[small]}\\
399: Large (5) & large & \Code{[LARGE]} \Code{[large]}\\
400: Vry Large (6) & very large & \Code{[VRY LARGE]} \Code{[vry large]}\\
401: Ext Large (7) & extra large & \Code{[EXT LARGE]} \Code{[ext large]}\\
402: \end{tabular}
403: \par\addvspace{4pt}
404: 
405: \Notes
406: \Topic{On the Screen}
407: The editing screen draws text of every size in the same cells as ordinary
408: text.  Press Reveal Codes to see which size text has.
409: 
410: \Topic{Size and Appearance Together}
411: Text can have a size and appearances at once; each is its own pair of codes
412: (see \InRef{Appearance}).
413: 
414: \SeeAlso{Appearance; Block, Attributes; Font}
415: \end{Feature}
416: 
417: \begin{Feature}{Status Line}
418: The status line is the bottom line of the InitechWord window.  While you
419: edit, it shows the name of the document at its left and where the cursor is
420: at its right.  InitechWord also uses it for menus, prompts, questions and
421: messages.
422: 
423: To read the status line,
424: \begin{WPSteps}
425: \item Look at the left end: the document's name once it has one, as
426: \Scr{A:\wbs MEMO}; or \Scr{Block on} while Block is on; or
427: \Scr{Move cursor; press Enter to retrieve.} while moved text waits.
428: \item Look at the right end: \Scr{Doc 1 Pg 1 Ln 1" Pos 1"}.
429: \end{WPSteps}
430: 
431: \BeginReplay{ref-status}
432: For example, type
433: \Type{Each line of a single-spaced paragraph is one sixth of an inch below the line before it, so Ln grows by 0.17 inch a line.}
434: The cursor is on the second line, after the period:
435: \WordScreen[rows=2]{ref-status}{Ln and Pos at the end of a two-line paragraph}
436: \EndReplay
437: 
438: \Notes
439: \Topic{Doc}
440: The document number is always 1.
441: 
442: \Topic{Ln}
443: Ln is the distance from the top edge of the page to the line the cursor is
444: on, in inches: \Scr{1"} for the first line under the one inch top margin,
445: then a sixth of an inch for each single-spaced line (\Scr{1.17"},
446: \Scr{1.33"}, \ldots).  With other line spacings, each line adds its own
447: spacing (see \InRef{Line Spacing}).
448: 
449: \Topic{Measurements}
450: Ln and Pos are shown in inches, to two decimal places at most, without
451: trailing zeros: \Scr{1"}, \Scr{1.5"}, \Scr{2.17"}.
452: 
453: \Topic{Pg}
454: Pg is the page the cursor is on (see \InRef{Page Break, Hard}).
455: 
456: \Topic{Pos}
457: Pos is the distance from the left edge of the paper to the cursor, in
458: inches: \Scr{1"} at the standard left margin, then a tenth of an inch for
459: each character.  While Caps Lock is on, Pos is shown as \Scr{POS} (see
460: \InRef{Caps Lock}).
461: 
462: \SeeAlso{Caps Lock; Cursor Movement; Line Spacing; Page Break, Hard}
463: \end{Feature}
